\documentclass[11pt,a4paper]{article}

\usepackage[T1]{fontenc}
\usepackage[utf8]{inputenc}
\usepackage[final]{microtype}
\usepackage[a4paper,margin=26mm,headheight=15pt]{geometry}
\usepackage{amsmath,amssymb,amsthm,mathtools}
\usepackage{newtxtext,newtxmath}
\usepackage{bm}
\usepackage{booktabs,longtable,array,tabularx}
\usepackage{enumitem}
\usepackage{xcolor}
\usepackage[most]{tcolorbox}
\usepackage{listings}
\usepackage{fancyhdr}
\usepackage{hyperref}
\usepackage[nameinlink,capitalise,noabbrev]{cleveref}

\definecolor{deepblue}{RGB}{25,63,112}
\definecolor{softblue}{RGB}{239,245,252}
\definecolor{softgray}{RGB}{247,247,247}
\definecolor{darkgray}{RGB}{72,72,72}
\definecolor{softgold}{RGB}{251,247,235}
\definecolor{deepgold}{RGB}{132,98,22}

\hypersetup{
  colorlinks=true,
  linkcolor=deepblue,
  citecolor=deepblue,
  urlcolor=deepblue,
  pdftitle={Asymptotic Inversion of the Gamma and Barnes G Functions},
  pdfauthor={Research note prepared for the ProveIt project},
  pdfsubject={Lambert-normalized transseries and power-logarithmic series reversal}
}

\pagestyle{fancy}
\fancyhf{}
\fancyhead[L]{\small Inverse Gamma and Barnes \(G\)}
\fancyhead[R]{\small Lambert-normalized transseries}
\fancyfoot[C]{\thepage}

\setlength{\parindent}{0pt}
\setlength{\parskip}{0.56em}
\setlist{nosep,leftmargin=2.0em}
\allowdisplaybreaks[3]
\numberwithin{equation}{section}

\newtheorem{theorem}{Theorem}[section]
\newtheorem{proposition}[theorem]{Proposition}
\newtheorem{lemma}[theorem]{Lemma}
\newtheorem{corollary}[theorem]{Corollary}
\theoremstyle{definition}
\newtheorem{definition}[theorem]{Definition}
\newtheorem{example}[theorem]{Example}
\theoremstyle{remark}
\newtheorem{remark}[theorem]{Remark}

\newtcolorbox{mainresultbox}[1][]{
  enhanced,
  breakable,
  colback=softblue,
  colframe=deepblue,
  boxrule=0.8pt,
  arc=1.2mm,
  left=3mm,right=3mm,top=2.5mm,bottom=2.5mm,
  title={#1},
  fonttitle=\bfseries
}
\newtcolorbox{methodbox}[1][]{
  enhanced,
  breakable,
  colback=softgray,
  colframe=darkgray,
  boxrule=0.5pt,
  arc=1mm,
  left=3mm,right=3mm,top=2mm,bottom=2mm,
  #1
}
\newtcolorbox{cautionbox}[1][]{
  enhanced,
  breakable,
  colback=softgold,
  colframe=deepgold,
  boxrule=0.6pt,
  arc=1mm,
  left=3mm,right=3mm,top=2mm,bottom=2mm,
  #1
}

\lstdefinestyle{pythonstyle}{
  language=Python,
  basicstyle=\ttfamily\footnotesize,
  columns=fullflexible,
  keepspaces=true,
  frame=single,
  framerule=0.4pt,
  breaklines=true,
  showstringspaces=false,
  keywordstyle=\bfseries,
  commentstyle=\itshape,
  xleftmargin=0.5em,
  xrightmargin=0.5em
}

\newcommand{\e}{\mathrm e}
\newcommand{\dd}{\,\mathrm d}
\newcommand{\Ord}{\mathcal O}
\newcommand{\littleo}{\mathrm o}
\newcommand{\cinv}{\langle-1\rangle}
\newcommand{\coeff}[2]{[\,#1^{#2}\,]}
\newcommand{\W}{W}
\newcommand{\Log}{\operatorname{Log}}
\newcommand{\sgn}{\operatorname{sgn}}
\newcommand{\Res}{\operatorname{Res}}
\newcommand{\Cfield}{\mathscr C}
\newcommand{\Tfield}{\mathscr T}
\newcommand{\val}{\nu}
\newcommand{\Bcal}{\mathcal B}
\newcommand{\Gcal}{\mathcal G}
\newcommand{\GamInv}{\Gamma_{+}^{\cinv}}
\newcommand{\BarnesInv}{G_{+}^{\cinv}}

\title{\textbf{Asymptotic Inversion of the Gamma and Barnes \(G\) Functions}\\[0.35em]
\Large Lambert-Normalized Transseries, Explicit Coefficients,\\
and a General Power--Logarithmic Reversion Calculus}
\author{Research note prepared for the ProveIt project}
\date{3 September 2026}

\begin{document}
\maketitle

\begin{abstract}
The large-value inverses of the Gamma function and the Barnes \(G\)-function
are governed by a common nonlinear scale: after taking a logarithm, their
dominant terms are respectively \(x\log x\) and \(\tfrac12x^2\log x\).
We develop a systematic inversion calculus for asymptotic transseries whose
leading block is
\[
  a x^p(\log x-\beta),
\]
or, more generally, \(a x^p(\log x-\beta)^\rho\).  The dominant block is
inverted exactly by a suitable branch of Lambert's \(W\)-function.  If
\(X\) denotes this exact leading inverse and \(q=X^{-1}\), the remaining
inverse is a triangular series
\[
  x=X+\sum_{n\geq0}C_n q^n
\]
with coefficients in a slow differential field generated by the Lambert
coordinate.  We prove formal existence and uniqueness, give a coefficient
recurrence and its first universal terms, derive a near-identity
Lagrange--Buer\-mann representation, describe Newton acceleration and
residual certification, and explain how algebraic and exponentially small
sectors are transported under inversion.

For the increasing real Gamma branch, with
\[
 T=\log\!\left(\frac{y}{\sqrt{2\pi}}\right),\qquad
 w=W_0(T/\e),\qquad X=T/w,\qquad s=w+1,
\]
we obtain
\[
 \GamInv(y)=\frac12+X+\frac{1}{24sX}
 -\frac{14s^2+10s+5}{5760s^3X^3}+\cdots,
\]
and display the expansion through \(X^{-9}\).  Centering the Stirling
series at \(x-\tfrac12\) forces every even inverse-power coefficient to
vanish.  We also give convergent large-\(y\) expansions for all inverse
branches tending to poles of \(\Gamma\), including the branch tending to
zero.

For the eventual positive Barnes branch, with
\[
 Y=\log y,\qquad w=W_0(4Y\e^{-3}),\qquad
 Z=2\sqrt{Y/w},\qquad D=w+1,
\]
we derive
\[
 \BarnesInv(y)=1+Z-\frac{\log(2\pi)}{D}
 +\frac{C_1(D)}{Z}+\frac{C_2(D)}{Z^2}
 +\frac{C_3(D)}{Z^3}+\cdots,
\]
with explicit rational functions \(C_1,C_2,C_3,C_4\).  Re-expansion of the
Lambert coordinate gives conventional nested-logarithm transseries, while
retaining \(W\) exactly preserves the correct sector ordering.  Exact
symbolic substitution and high-precision numerical tests verify all
reported coefficients.
\end{abstract}

\begin{mainresultbox}[Principal real branches at a glance]
Write compositional inverses with the symbol \(f^{\cinv}\), so that they
cannot be confused with reciprocals.

For \(y\to+\infty\), the increasing real Gamma branch is
\begin{align*}
 \GamInv(y)={}&\frac12+X+\frac{1}{24sX}
 -\frac{14s^2+10s+5}{5760s^3X^3}\\
 &+\frac{2232s^4+1176s^3+714s^2+350s+105}
 {2903040s^5X^5}\\
 &-\frac{822960s^6+292536s^5+136956s^4+68040s^3
 +32970s^2+12250s+2625}
 {1393459200s^7X^7}\\
 &+\frac{309052800s^8+77920128s^7+26024592s^6+10019592s^5
 +4516028s^4}{367873228800s^9X^9}\\
 &\quad+\frac{2023560s^3+812350s^2+242550s+40425}
 {367873228800s^9X^9}
 +\Ord\!\left(\frac{1}{sX^{11}}\right),
\end{align*}
where
\[
 T=\log(y/\sqrt{2\pi}),\qquad w=W_0(T/\e),\qquad X=T/w,
 \qquad s=w+1.
\]

For the eventual positive Barnes branch, set
\[
 Y=\log y,\qquad w=W_0(4Y\e^{-3}),\qquad
 Z=2\sqrt{Y/w},\qquad D=w+1,
\]
and put \(\lambda=\log(2\pi)\), \(\alpha=\log A\), where \(A\) is
Glaisher's constant.  Then
\[
 \BarnesInv(y)=1+Z+C_0+\frac{C_1}{Z}+\frac{C_2}{Z^2}
 +\frac{C_3}{Z^3}+\Ord(Z^{-4}),
\]
with
\begin{align*}
 C_0={}&-\frac{\lambda}{D},\\
 C_1={}&\frac1{12}+\frac{2\alpha}{D}
 +\frac{\lambda^2}{2D^2}-\frac{\lambda^2}{D^3},\\
 C_2={}&\frac{4\lambda\alpha}{D^3}
 +\frac{4\lambda^3}{3D^4}-\frac{2\lambda^3}{D^5},
\end{align*}
and \(C_3\) is displayed in \cref{thm:barnes-main}.
\end{mainresultbox}

\tableofcontents
\newpage

\section{Introduction}

The Gamma and Barnes \(G\)-functions grow faster than exponentials but more
slowly than iterated exponentials.  Their logarithms have the leading
forms
\begin{align}
 \log\Gamma(x)&=x\log x-x+\text{lower blocks},\label{eq:intro-gamma}\\
 \log G(x+1)&=\frac12x^2\log x-\frac34x^2+\text{lower blocks}.
 \label{eq:intro-barnes}
\end{align}
Consequently, their large-value inverse functions are not described by a
plain Puiseux or inverse-power series.  The first inversion already
requires Lambert's \(W\)-function; after this exact normalization, the
remaining corrections form a power series in the reciprocal leading
inverse, with coefficients that vary only logarithmically.

The leading Lambert approximation to the principal inverse Gamma function
is classical and is discussed, for example, by Borwein and Corless and by
Jeffrey and Watt~\cite{BorweinCorless,JeffreyWatt}.  The analytic branch
structure of inverse Gamma functions was studied by Uchiyama and
Pedersen~\cite{Uchiyama,Pedersen}.  Our purpose here is different: we carry
the large-argument inversion to arbitrary asymptotic order, make the
coefficient recursion explicit, apply the same machinery to Barnes \(G\),
and isolate a reusable theory for power--logarithmic transseries.

The linked ProveIt drafts develop a broad calculus for logarithmic
transseries and near-identity series reversal~\cite{ProveItTransseries}.
The present article uses the same organizing principle but adapts the
coordinate system to maps whose leading derivative is not asymptotic to
one.  The decisive step is to invert the dominant power--logarithmic block
\emph{exactly}.  In the resulting Lambert coordinate, the correction
problem becomes near-identity and triangular.

\subsection{What is meant by a transseries here}

Three nested levels occur.

\begin{enumerate}
\item The \emph{Lambert-normalized outer expansion} is a series in
      \(q=X^{-1}\), where \(X\) is the exact inverse of the dominant block.
      Its coefficients are rational functions, or more general slow
      transseries, in a Lambert coordinate \(w\sim\log Y\).
\item Expanding \(W\) itself in \(\log Y\) and \(\log\log Y\) produces a
      conventional logarithmic transseries.
\item The Stirling series are factorially divergent.  Exponential
      improvement adds beyond-all-orders sectors; inversion transports
      each such sector by division by the derivative of the algebraic
      part, followed by triangular nonlinear coupling.
\end{enumerate}

The first level is normally the best representation for computation and
for clean asymptotic bookkeeping.  The second is useful when an answer is
required entirely in elementary iterated logarithms.  The third records
the optimal-truncation and Stokes structure.

\subsection{Branch conventions}

On the positive real axis, \(\Gamma\) decreases from \(+\infty\) at
\(0^+\) to its unique positive minimum and then increases to
\(+\infty\).  Thus a sufficiently large positive value has two positive
preimages.  We write
\begin{itemize}
\item \(\Gamma_-^{\cinv}(y)\to0^+\) for the left branch;
\item \(\Gamma_+^{\cinv}(y)\to+\infty\) for the increasing, or principal,
      branch.
\end{itemize}
In addition, there are inverse branches approaching every pole
\(-n\), \(n=0,1,2,\ldots\).  These pole branches have convergent expansions
in \(1/y\), unlike the Stirling branch at \(+\infty\).

The Barnes \(G\)-function is not globally one-to-one on the positive axis,
but \((\log G)'(x)\sim x\log x\); hence it is strictly increasing for all
sufficiently large \(x\).  We denote its eventual positive inverse by
\(G_+^{\cinv}\).

For complex inversion, one first chooses a branch of \(\Log y\), then a
branch \(W_k\), and finally a sector in which the chosen Stirling expansion
and logarithm are consistent.  The coefficient formulae below are
unchanged; only the slow coordinates and branch labels change.

\subsection{Summary of the method}

The reusable pipeline is:
\begin{enumerate}
\item take a logarithm of the target special function when appropriate;
\item translate the argument and subtract constants to put the forward
      expansion in a sparse normal form;
\item invert the leading block by Lambert \(W\);
\item set the true inverse equal to the leading inverse plus an additive
      correction;
\item expand the normalized residual in \(q=X^{-1}\) and solve one linear
      equation per coefficient;
\item certify the truncation by substitution or by one Newton step.
\end{enumerate}

\section{Forward asymptotics and normal forms}

\subsection{Gamma in the half-shifted coordinate}

The standard Stirling expansion is~\cite{DLMFGamma,Olver}
\begin{equation}
 \log\Gamma(x)\sim
 \left(x-\frac12\right)\log x-x+\frac12\log(2\pi)
 +\sum_{k\ge1}\frac{B_{2k}}{2k(2k-1)x^{2k-1}}.
 \label{eq:stirling-standard}
\end{equation}
For inversion, the translated coordinate
\begin{equation}
 z=x-\frac12
 \label{eq:gamma-center}
\end{equation}
is substantially better.  The Bernoulli-polynomial form of Stirling's
series gives
\begin{equation}
 \log\Gamma\!\left(z+\frac12\right)
 \sim z(\log z-1)+c_\Gamma
 +\sum_{k\ge1}\frac{a_k}{z^{2k-1}},
 \qquad c_\Gamma=\frac12\log(2\pi),
 \label{eq:gamma-centered}
\end{equation}
where
\begin{equation}
 a_k=\frac{B_{2k}(1/2)}{2k(2k-1)}
 =\frac{(2^{1-2k}-1)B_{2k}}{2k(2k-1)}.
 \label{eq:ak-def}
\end{equation}
The first values are
\begin{equation}
 a_1=-\frac1{24},\quad
 a_2=\frac7{2880},\quad
 a_3=-\frac{31}{40320},\quad
 a_4=\frac{127}{215040},\quad
 a_5=-\frac{511}{608256}.
 \label{eq:ak-first}
\end{equation}
Only odd inverse powers occur.  This sparsity propagates to the inverse
series and is the main reason for centering at one half.

\subsection{Barnes G in the natural shifted coordinate}

Barnes' function satisfies
\begin{equation}
 G(z+1)=\Gamma(z)G(z),\qquad G(1)=1.
 \label{eq:barnes-recurrence}
\end{equation}
For \(|\arg z|\leq\pi-\delta\), its logarithm has the asymptotic
expansion~\cite{DLMFBarnes,Nemes}
\begin{align}
 \log G(z+1)\sim{}&
 \frac12z^2\log z-\frac34z^2
 +\frac{\lambda}{2}z-\frac1{12}\log z+\kappa
 +\sum_{k\ge1}\frac{b_k}{z^{2k}},
 \label{eq:barnes-forward}\\
 \lambda={}&\log(2\pi),\qquad
 \kappa=\zeta'(-1)=\frac1{12}-\alpha,
 \qquad \alpha=\log A,\label{eq:barnes-constants}\\
 b_k={}&\frac{B_{2k+2}}{2k(2k+2)}.
 \label{eq:bk-def}
\end{align}
Here \(A\) is Glaisher's constant.  Thus
\begin{equation}
 b_1=-\frac1{240},\qquad b_2=\frac1{1008},\qquad
 b_3=-\frac1{1440},\quad\ldots
 \label{eq:bk-first}
\end{equation}
The displayed form follows either directly from a standard Barnes
expansion or by inserting Stirling's series into the DLMF identity for
\(\log G(z+1)\).  Nemes gives explicit remainder bounds and an
exponentially improved form~\cite{Nemes}.

\subsection{Normal-form principle}

A translation \(x=z+\theta\) and a target shift \(Y\mapsto Y-c\) do not
change the inversion problem, but they can remove entire slow blocks.
The Gamma translation \(\theta=\tfrac12\) changes the standard Stirling
series into \cref{eq:gamma-centered}, eliminating a logarithmic block at
outer order zero and revealing an odd parity.  For Barnes \(G\), the
natural recurrence already singles out \(z=x-1\); the linear term
\(\tfrac\lambda2z\) is retained and generates the first inverse correction.

\begin{methodbox}[title={Practical normalization rule}]
Before reversing a power--logarithmic expansion, choose translations and
constant target shifts to maximize the valuation of the remainder.  A
single removed low-order block can simplify every subsequent coefficient.
\end{methodbox}

\section{Exact Lambert normalization of a power--logarithmic leader}

\subsection{The basic leader}

Let
\begin{equation}
 F_0(x)=a x^p(\log x-\beta),
 \qquad a\neq0,\quad p\neq0.
 \label{eq:F0-basic}
\end{equation}
To solve \(F_0(X)=Y\), set
\begin{equation}
 w=W_k\!\left(\frac{pY}{a}\e^{-p\beta}\right),
 \qquad
 X=\exp\!\left(\beta+\frac{w}{p}\right)
 =\left(\frac{pY}{aw}\right)^{1/p}.
 \label{eq:Lambert-normalizer}
\end{equation}
Then
\begin{equation}
 p(\log X-\beta)=w,
 \qquad
 aX^p\frac{w}{p}=Y.
 \label{eq:Lambert-identities}
\end{equation}
On the positive real axis with \(a,p,Y>0\), the relevant branch is
\(W_0\).

\begin{proposition}[Exact dominant inversion]\label{prop:dominant-inversion}
The value \(X\) in \cref{eq:Lambert-normalizer} solves
\(F_0(X)=Y\) exactly on every compatible choice of logarithm, root, and
Lambert branch.
\end{proposition}

\begin{proof}
Put \(u=p(\log X-\beta)\).  Then
\[
 F_0(X)=\frac{a}{p}\e^{p\beta}u\e^u.
\]
Thus \(F_0(X)=Y\) is equivalent to
\(u\e^u=(pY/a)\e^{-p\beta}\), so \(u=W_k((pY/a)\e^{-p\beta})\).
Exponentiating \(\log X=\beta+u/p\) proves the claim.
\end{proof}

The derivative at the leading inverse is
\begin{equation}
 F_0'(X)=aX^{p-1}(w+1).
 \label{eq:F0prime}
\end{equation}
The ubiquitous denominator \(w+1\) in inverse coefficients is therefore
not accidental: it is the normalized Jacobian of the dominant map.

\subsection{A general logarithmic power}

The same exact reduction works for
\begin{equation}
 F_{0,\rho}(x)=a x^p(\log x-\beta)^\rho,
 \qquad \rho\neq0.
 \label{eq:F0-rho}
\end{equation}
Choose compatible branches and define
\begin{equation}
 w=W_k\!\left(
 \frac{p}{\rho}
 \left(\frac{Y}{a\e^{p\beta}}\right)^{1/\rho}
 \right),
 \qquad
 X=\exp\!\left(\beta+\frac{\rho}{p}w\right).
 \label{eq:rho-normalizer}
\end{equation}
Then \(F_{0,\rho}(X)=Y\).  The normalized linear coefficient in an
additive correction is proportional to
\begin{equation}
 p+\frac{\rho}{\log X-\beta}=p\left(1+\frac1w\right),
 \label{eq:rho-jacobian}
\end{equation}
so the same triangular inversion mechanism applies away from \(w=-1\).
The remainder of the article develops explicit universal formulas for the
case \(\rho=1\), which includes both special functions of interest.

\section{A formal reversion calculus in the Lambert coordinate}

\subsection{The transseries class}

Consider a forward expansion
\begin{equation}
 F(x)\sim a x^p(\log x-\beta)
 +\sum_{m\ge1}x^{p-m}P_m(\log x),
 \label{eq:general-forward}
\end{equation}
where the \(P_m\) may be polynomials, rational functions, or elements of a
slow differential field containing iterated logarithms.  Put
\begin{equation}
 R_m(s)=\frac{P_m(s)}a,
 \qquad
 s=\log X=\beta+\frac{w}{p},
 \qquad q=X^{-1},
 \label{eq:Rmsq}
\end{equation}
and seek
\begin{equation}
 x=X+U(q,w),
 \qquad
 U(q,w)=\sum_{n\ge0}C_n(w)q^n.
 \label{eq:U-series}
\end{equation}
The slow variable \(w=p(s-\beta)\) grows logarithmically while \(q\) is the
primary small parameter.  Formal series are taken in the complete
\(q\)-adic ring
\begin{equation}
 \Tfield=\Cfield(w)[[q]],
 \label{eq:Tfield}
\end{equation}
where \(\Cfield(w)\) is closed under the derivatives and Taylor
translations required below.

\subsection{The normalized residual equation}

Since \(x=X(1+qU)\), direct substitution into \cref{eq:general-forward}
leads to the normalized equation
\begin{align}
 \mathscr E(q,U;w)={}&
 \frac{(1+qU)^p\left(\dfrac wp+\log(1+qU)\right)-\dfrac wp}{q}
 \notag\\
 &+\sum_{m\ge1}q^{m-1}(1+qU)^{p-m}
 R_m\!\left(s+\log(1+qU)\right)=0.
 \label{eq:normalized-E}
\end{align}
Indeed,
\begin{equation}
 F(X+U)-Y=aX^{p-1}\mathscr E(q,U;w).
 \label{eq:residual-scale}
\end{equation}
Expanding \(\mathscr E\) at \(q=0\) gives
\begin{equation}
 \mathscr E(q,U;w)=(w+1)U+R_1(s)+\Ord(q).
 \label{eq:E-leading}
\end{equation}
Thus every new coefficient enters linearly with multiplier \(w+1\).

\begin{theorem}[Lambert-normalized formal inverse]\label{thm:formal-inverse}
Assume that \(w+1\) is invertible in the slow coefficient field.  Then
\cref{eq:normalized-E} has a unique solution
\(U\in\Cfield(w)[[q]]\).  Its coefficients are determined triangularly by
\begin{equation}
 C_n(w)=-\frac1{w+1}
 \coeff qn{\mathscr E\!\left(q,U_{<n};w\right)},
 \qquad
 U_{<n}=\sum_{j=0}^{n-1}C_j(w)q^j.
 \label{eq:triangular-recurrence}
\end{equation}
If \(P_1=\cdots=P_{m_0-1}=0\), then
\(C_0=\cdots=C_{m_0-2}=0\).
\end{theorem}

\begin{proof}
Suppose \(C_0,\ldots,C_{n-1}\) are known and insert
\(U=U_{<n}+C_nq^n+\Ord(q^{n+1})\) into \cref{eq:normalized-E}.  The
coefficient of \(q^n\) contributed by the first derivative in \(U\) is
\((w+1)C_n\).  Every nonlinear occurrence of \(C_nq^n\) carries at least
one additional factor of \(q\), and every lower block in the second line
of \cref{eq:normalized-E} does the same.  Hence no other term at order
\(q^n\) contains \(C_n\).  Solving the resulting linear equation gives
\cref{eq:triangular-recurrence}.  Induction proves existence and
uniqueness.  The valuation statement follows because the first nonzero
forward block \(P_{m_0}\) enters at order \(q^{m_0-1}\).
\end{proof}

\subsection{Closed coefficient form}

For an explicit partition formula, define
\begin{align}
 f_j(w)&=\coeff tj{(1+t)^p\left(\frac wp+\log(1+t)\right)},
 \label{eq:fj-def}\\
 \mathcal R_{m,j}(s)&=\coeff tj{(1+t)^{p-m}
 R_m\!\left(s+\log(1+t)\right)}.
 \label{eq:Rmj-def}
\end{align}
Then \(f_1=w+1\), and
\begin{equation}
 \mathscr E(q,U;w)=
 \sum_{j\ge1}f_j(w)q^{j-1}U^j
 +\sum_{m\ge1}\sum_{j\ge0}\mathcal R_{m,j}(s)
 q^{m+j-1}U^j.
 \label{eq:E-double-series}
\end{equation}
Consequently,
\begin{align}
 (w+1)C_n={}&-
 \sum_{j=2}^{n+1}f_j(w)
 \coeff q{n-j+1}{U_{<n}^j}
 \notag\\
 &-\sum_{m=1}^{n+1}\sum_{j=0}^{n-m+1}
 \mathcal R_{m,j}(s)
 \coeff q{n-m-j+1}{U_{<n}^j}.
 \label{eq:partition-recurrence}
\end{align}
Only finitely many terms contribute at each \(n\).  Formula
\cref{eq:partition-recurrence} is convenient for symbolic implementation
with Bell polynomials or ordinary truncated-series arithmetic.

\subsection{The first three universal coefficients}

Write
\begin{equation}
 d=w+1,
 \qquad r_m=R_m(s),
 \qquad \dot r_m=\frac{\dd r_m}{\dd s}.
 \label{eq:dots}
\end{equation}
Define
\begin{align}
 \Lambda_2&=\frac{(p-1)w+2p-1}{2},
 \label{eq:Lambda2}\\
 \Lambda_3&=\frac{(p-1)(p-2)w+3p^2-6p+2}{6},
 \label{eq:Lambda3}\\
 B_1&=(p-1)r_1+\dot r_1,
 \label{eq:B1}\\
 B_2&=\binom{p-1}{2}r_1
 +\left(p-\frac32\right)\dot r_1+\frac12\ddot r_1,
 \label{eq:B2}\\
 D_2&=(p-2)r_2+\dot r_2.
 \label{eq:D2}
\end{align}
Then
\begin{align}
 C_0={}&-\frac{r_1}{d},
 \label{eq:C0-univ}\\
 C_1={}&-\frac{\Lambda_2C_0^2+B_1C_0+r_2}{d},
 \label{eq:C1-univ}\\
 C_2={}&-\frac{2\Lambda_2C_0C_1+\Lambda_3C_0^3
 +B_1C_1+B_2C_0^2+D_2C_0+r_3}{d}.
 \label{eq:C2-univ}
\end{align}
These formulas already recover the first three Barnes coefficients and
explain the common denominator structure.

\subsection{Near-identity reduction and Lagrange--Buermann inversion}

There is a second, coordinate-free derivation.  Let
\begin{equation}
 \Psi(x)=F_0^{\cinv}(F(x))=x+H(x).
 \label{eq:Psi-near-id}
\end{equation}
If \(F(x)=Y\) and \(X=F_0^{\cinv}(Y)\), then \(\Psi(x)=X\), so
\begin{equation}
 x=\Psi^{\cinv}(X).
 \label{eq:Psi-inverse}
\end{equation}
Moreover,
\begin{equation}
 H(x)\sim\frac{F(x)-F_0(x)}{F_0'(x)}+\cdots,
 \label{eq:H-leading}
\end{equation}
which belongs to the slow \(q\)-adic algebra.  Differentiation with respect
to \(X\) raises the primary \(q\)-valuation because
\begin{equation}
 \frac{\dd}{\dd X}\bigl(q^nA(w)\bigr)
 =q^{n+1}\bigl(pA'(w)-nA(w)\bigr).
 \label{eq:derivative-valuation}
\end{equation}
The usual Lagrange--Buermann formula is therefore summable blockwise.

\begin{theorem}[Lagrange representation at infinity]\label{thm:LB-infinity}
Under the hypotheses of \cref{thm:formal-inverse},
\begin{equation}
 F^{\cinv}(Y)=X+
 \sum_{n\ge1}\frac{(-1)^n}{n!}
 \left(\frac{\dd}{\dd X}\right)^{n-1}H(X)^n,
 \qquad H=F_0^{\cinv}\!\circ F-\operatorname{id}.
 \label{eq:LB-infinity}
\end{equation}
At every fixed power of \(q\), only finitely many summands contribute.
\end{theorem}

\begin{proof}
Apply the formal Lagrange--Buermann formula to the near-identity map
\(X\mapsto X+H(X)\).  By \cref{eq:derivative-valuation}, the \(n\)-th
summand has primary valuation at least \(n-1\).  Hence the family is
summable in the \(q\)-adic topology, and the resulting series is the unique
inverse tangent to the identity.
\end{proof}

The direct residual recurrence \cref{eq:triangular-recurrence} is usually
more efficient, while \cref{eq:LB-infinity} clarifies the relation to
classical series reversion and provides an independent check.

\subsection{Newton acceleration in the formal ring}

Starting with any truncation \(U^{[0]}\), define
\begin{equation}
 U^{[j+1]}=U^{[j]}-
 \frac{\mathscr E(q,U^{[j]};w)}
 {\partial_U\mathscr E(q,U^{[j]};w)}.
 \label{eq:formal-Newton}
\end{equation}
Since \(\partial_U\mathscr E=(w+1)+\Ord(q)\) is invertible, Newton iteration
approximately doubles the number of correct \(q\)-coefficients at each
step.  This is preferable to coefficient-by-coefficient recursion at very
high order.

\section{From formal series to asymptotic and exponential sectors}

\subsection{Residual certification}

The normalized equation gives a particularly simple error test.  Suppose a
truncation \(U_N\) satisfies
\begin{equation}
 \mathscr E(q,U_N;w)=\Ord(q^{N+1}M(w))
 \label{eq:E-residual-order}
\end{equation}
for a slow coefficient \(M\).  Then
\begin{equation}
 F(X+U_N)-Y
 =\Ord\!\left(aX^{p-N-2}M(w)\right).
 \label{eq:forward-residual-order}
\end{equation}
If \(F\) is real, eventually increasing, and
\begin{equation}
 F'(x)\ge cX^{p-1}|w+1|
 \label{eq:derivative-lower}
\end{equation}
through the interval joining the exact and approximate roots, the mean
value theorem gives
\begin{equation}
 F^{\cinv}(Y)-(X+U_N)
 =\Ord\!\left(\frac{q^{N+1}M(w)}{w+1}\right).
 \label{eq:inverse-error-general}
\end{equation}
This turns symbolic cancellation of the residual into an analytic error
estimate once a differentiated forward asymptotic expansion is available.

For numerical work one can use the a posteriori bound
\begin{equation}
 |x-\widehat x|\le
 \frac{|F(\widehat x)-Y|}{\inf_{I}|F'|},
 \label{eq:apost-bound}
\end{equation}
where \(I\) joins \(x\) and \(\widehat x\).  A single ordinary Newton step
then typically squares the already small transseries error.

\subsection{Transport of genuinely exponential sectors}

Suppose an exponentially improved forward transseries has the schematic
form
\begin{equation}
 F(x;\bm\sigma)=F_{\rm alg}(x)
 +\sum_{\bm m\ne0}\bm\sigma^{\bm m}R_{\bm m}(x),
 \label{eq:forward-exp-sectors}
\end{equation}
where \(R_{\bm m}\) contains an exponential monomial such as
\(\exp[-\bm m\cdot\bm\omega\,x^\nu]\).  Let
\(x_{\rm alg}=F_{\rm alg}^{\cinv}(Y)\) and seek
\begin{equation}
 x=x_{\rm alg}+\sum_{\bm m\ne0}
 \bm\sigma^{\bm m}\Delta_{\bm m}(x_{\rm alg}).
 \label{eq:inverse-exp-sectors}
\end{equation}
The first sector is universal:
\begin{equation}
 \Delta_{\bm e_j}(x_{\rm alg})
 =-\frac{R_{\bm e_j}(x_{\rm alg})}
 {F_{\rm alg}'(x_{\rm alg})}.
 \label{eq:first-exp-sector}
\end{equation}
Higher sectors are obtained recursively by expanding
\(F(x_{\rm alg}+\Delta;\bm\sigma)-Y\) and extracting one exponential
weight at a time.  They contain derivatives of lower sectors and the usual
Bell-polynomial couplings.

\begin{theorem}[Transseries implicit inversion]\label{thm:transseries-implicit}
In a well-based transseries field, if \(F_{\rm alg}'(x_{\rm alg})\) is
invertible, the sector equations generated by
\cref{eq:forward-exp-sectors,eq:inverse-exp-sectors} are triangular in
exponential weight.  They have a unique solution tangent to
\(x_{\rm alg}\), and their first equation is \cref{eq:first-exp-sector}.
\end{theorem}

\begin{proof}
Expand each forward sector at \(x_{\rm alg}+\Delta\).  At a fixed nonzero
exponential weight \(\bm m\), the only term containing the new unknown
\(\Delta_{\bm m}\) linearly is
\(F_{\rm alg}'(x_{\rm alg})\Delta_{\bm m}\).  Every other contribution is
a product or derivative of lower-weight sectors.  Well-basedness ensures
that only finitely many decompositions contribute to a prescribed weight.
Division by the invertible derivative gives the unique recursion.
\end{proof}

The Bernoulli coefficients in both Stirling series grow factorially.
Least-term truncation occurs at an index proportional to \(2\pi X\) or
\(2\pi Z\); hence the remaining forward error has exponential order
\(\exp[-2\pi X+\Ord(\log X)]\) or
\(\exp[-2\pi Z+\Ord(\log Z)]\).  The corresponding inverse scales are,
schematically,
\begin{align}
 \delta x_\Gamma&=
 \frac{\exp[-2\pi X+\Ord(\log X)]}{\log X},
 \label{eq:gamma-exp-scale}\\
 \delta x_G&=
 \frac{\exp[-2\pi Z+\Ord(\log Z)]}{Z\log Z}.
 \label{eq:barnes-exp-scale}
\end{align}
Precise Stokes multipliers are inherited from the exponentially improved
forward expansions.  For Barnes \(G\), such expansions and Berry smoothing
are developed by Nemes~\cite{Nemes}; analogous exponential improvement for
Stirling's series is classical~\cite{DLMFGamma,ParisKaminski}.

\section{The increasing inverse Gamma function}

\subsection{Lambert coordinates}

Let \(y\to+\infty\), set
\begin{equation}
 Y=\log y,
 \qquad T=Y-c_\Gamma=\log\!\left(\frac{y}{\sqrt{2\pi}}\right),
 \label{eq:gamma-T}
\end{equation}
and define
\begin{equation}
 w=W_0(T/\e),
 \qquad X=\e^{w+1}=\frac{T}{w},
 \qquad s=w+1=\log X,
 \qquad q=X^{-1}.
 \label{eq:gamma-coordinates}
\end{equation}
Then
\begin{equation}
 X(\log X-1)=T.
 \label{eq:gamma-leading-exact}
\end{equation}
Writing
\begin{equation}
 \Gamma_+^{\cinv}(y)=\frac12+X+U_\Gamma(q,s)
 \label{eq:gamma-U-def}
\end{equation}
and using \cref{eq:gamma-centered} gives
\begin{align}
 \mathscr E_\Gamma(q,U;s)={}&
 \frac{(1+qU)\bigl(s-1+\log(1+qU)\bigr)-(s-1)}{q}
 \notag\\
 &+\sum_{k\ge1}a_kq^{2k-1}(1+qU)^{-(2k-1)}=0.
 \label{eq:E-gamma}
\end{align}
Its linear coefficient is \(s\).

\subsection{Parity and all-orders recurrence}

\begin{proposition}[Odd inverse-power parity]\label{prop:gamma-parity}
The unique formal solution of \cref{eq:E-gamma} satisfies
\begin{equation}
 U_\Gamma(-q,s)=-U_\Gamma(q,s).
 \label{eq:gamma-odd}
\end{equation}
Consequently,
\begin{equation}
 U_\Gamma(q,s)=\sum_{n\ge1}A_{2n-1}(s)q^{2n-1}.
 \label{eq:gamma-odd-series}
\end{equation}
\end{proposition}

\begin{proof}
A direct check gives
\(\mathscr E_\Gamma(-q,-U;s)=-\mathscr E_\Gamma(q,U;s)\).
Therefore, if \(U(q,s)\) is a solution, so is \(-U(-q,s)\).  Uniqueness
from \cref{thm:formal-inverse} makes them equal.
\end{proof}

The coefficient recursion is
\begin{equation}
 A_n(s)=-\frac1s\coeff qn{
 \mathscr E_\Gamma\!\left(q,\sum_{j=0}^{n-1}A_j(s)q^j;s\right)},
 \qquad A_{2j}=0.
 \label{eq:gamma-recurrence}
\end{equation}
At large \(s\), the direct contribution of the next Stirling coefficient
dominates:
\begin{equation}
 A_{2n-1}(s)=-\frac{a_n}{s}+\Ord(s^{-2}).
 \label{eq:gamma-leading-An}
\end{equation}
Because \(a_n\) has factorial growth, the outer Gamma inverse series is an
asymptotic, generally divergent, transseries.

\subsection{Explicit expansion}

\begin{theorem}[Right inverse Gamma through five correction sectors]
\label{thm:gamma-main}
With \(T,w,X,s\) as in \cref{eq:gamma-T,eq:gamma-coordinates},
\begin{align}
 \GamInv(y)={}&\frac12+X
 +\frac{1}{24sX}
 -\frac{14s^2+10s+5}{5760s^3X^3}
 \notag\\
 &+\frac{2232s^4+1176s^3+714s^2+350s+105}
 {2903040s^5X^5}
 \notag\\
 &-\frac{822960s^6+292536s^5+136956s^4+68040s^3
 +32970s^2+12250s+2625}
 {1393459200s^7X^7}
 \notag\\
 &+\frac{309052800s^8+77920128s^7+26024592s^6+10019592s^5
 +4516028s^4}{367873228800s^9X^9}
 \notag\\
 &\quad+\frac{2023560s^3+812350s^2+242550s+40425}
 {367873228800s^9X^9}
 +\Ord\!\left(\frac{1}{sX^{11}}\right).
 \label{eq:gamma-main-expansion}
\end{align}
The expansion is Poincare-asymptotic on the increasing real branch.  The
same coefficient functions apply on compatible complex branches after
replacing \(W_0\) and the logarithms by the selected branches.
\end{theorem}

\begin{proof}
Insert \(U=\sum A_{2n-1}q^{2n-1}\) into
\cref{eq:E-gamma}.  Applying \cref{eq:gamma-recurrence} successively gives
\begin{align*}
 A_1={}&\frac1{24s},\\
 A_3={}&-\frac{14s^2+10s+5}{5760s^3},\\
 A_5={}&\frac{2232s^4+1176s^3+714s^2+350s+105}{2903040s^5},\\
 A_7={}&-\frac{822960s^6+292536s^5+136956s^4+68040s^3
 +32970s^2+12250s+2625}{1393459200s^7},\\
 A_9={}&\frac{309052800s^8+77920128s^7+26024592s^6+10019592s^5
 +4516028s^4}{367873228800s^9}\\
 &+\frac{2023560s^3+812350s^2+242550s+40425}{367873228800s^9}.
\end{align*}
Substitution cancels \(\mathscr E_\Gamma\) through order \(q^9\); the next
possible term has order \(q^{11}\).  The differentiated Stirling remainder
and \cref{eq:inverse-error-general} promote the formal identity to the
stated real asymptotic expansion.
\end{proof}

\subsection{Derivative and residual checks}

If \(x=\GamInv(y)\), then
\begin{equation}
 \frac{\dd x}{\dd Y}=\frac1{\psi(x)},
 \qquad
 \frac{\dd x}{\dd y}=\frac1{y\psi(x)},
 \label{eq:gamma-inverse-derivative}
\end{equation}
where \(\psi\) is the digamma function.  Differentiating a transseries
truncation and comparing with \cref{eq:gamma-inverse-derivative} is an
independent coefficient check.  For numerical refinement, use
\begin{equation}
 x_{\rm new}=x-\frac{\log\Gamma(x)-Y}{\psi(x)}.
 \label{eq:gamma-Newton}
\end{equation}

\section{Pole branches of the inverse Gamma function}

The Lambert sector describes the branch tending to \(+\infty\).  Every
pole of \(\Gamma\) gives a different, locally convergent inverse expansion.
Let
\begin{equation}
 x=-n+\varepsilon,
 \qquad n=0,1,2,\ldots,
 \qquad r=\frac1y.
 \label{eq:pole-variables}
\end{equation}
The reciprocal Gamma function is entire and has a simple zero at \(-n\):
\begin{align}
 \frac1{\Gamma(-n+\varepsilon)}
 & =(-1)^n n!\,\varepsilon\,\Phi_n(\varepsilon),
 \label{eq:recip-gamma-pole}\\
 \Phi_n(\varepsilon)
 & =\prod_{k=1}^n\left(1-\frac{\varepsilon}{k}\right)
 \frac1{\Gamma(1+\varepsilon)},
 \qquad \Phi_n(0)=1.
 \label{eq:Phi-n}
\end{align}
Thus \(r=(-1)^nn!\varepsilon\Phi_n(\varepsilon)\), and ordinary Lagrange
inversion gives a convergent power series.

\begin{theorem}[All pole branches]\label{thm:gamma-pole-branches}
For sufficiently small \(|r|\), the inverse branch tending to \(-n\) is
\begin{equation}
 \Gamma_{[-n]}^{\cinv}(y)
 =-n+\sum_{m\ge1}\frac{t^m}{m}
 \coeff{\varepsilon}{m-1}{\Phi_n(\varepsilon)^{-m}},
 \qquad t=\frac{(-1)^n}{n!y}.
 \label{eq:pole-Lagrange}
\end{equation}
If \(\psi_n=\psi(n+1)=H_n-\gamma\), then
\begin{align}
 \Gamma_{[-n]}^{\cinv}(y)={}&-n+t+\psi_n t^2
 +\frac{3\psi_n^2+H_n^{(2)}+\zeta(2)}{2}t^3
 +\Ord(t^4).
 \label{eq:pole-first-terms}
\end{align}
For positive \(y\), the leading displacement from the pole has sign
\((-1)^n\).
\end{theorem}

\begin{proof}
Equation \cref{eq:recip-gamma-pole} follows from
\[
 \Gamma(-n+\varepsilon)=
 \frac{\Gamma(1+\varepsilon)}
 {\varepsilon(\varepsilon-1)\cdots(\varepsilon-n)}.
\]
Apply the classical Lagrange formula to
\(t=\varepsilon\Phi_n(\varepsilon)\).  Expanding
\[
 \log\Phi_n(\varepsilon)
 =-\psi_n\varepsilon
 -\frac{H_n^{(2)}+\zeta(2)}2\varepsilon^2
 +\frac{\zeta(3)-H_n^{(3)}}3\varepsilon^3+\cdots
\]
gives the displayed coefficients.
\end{proof}

\subsection{The branch tending to zero}

For \(n=0\), \(t=r=1/y\).  Expanding further yields
\begin{align}
 \Gamma_-^{\cinv}(y)={}&r-\gamma r^2
 +\frac{18\gamma^2+\pi^2}{12}r^3
 -\frac{8\gamma^3+\pi^2\gamma+\zeta(3)}3r^4
 \notag\\
 &+\frac{7500\gamma^4+1500\pi^2\gamma^2
 +2400\gamma\zeta(3)+29\pi^4}{1440}r^5
 \notag\\
 &-\frac{1296\gamma^5+360\pi^2\gamma^3
 +720\gamma^2\zeta(3)+17\pi^4\gamma}{120}r^6
 \notag\\
 &-\frac{24\zeta(5)+20\pi^2\zeta(3)}{120}r^6
 +\Ord(r^7).
 \label{eq:gamma-left-expansion}
\end{align}
Unlike \cref{eq:gamma-main-expansion}, this is a genuinely convergent local
series because it inverts an analytic map with nonzero derivative at the
origin.

\section{The eventual positive inverse of Barnes G}

\subsection{Lambert coordinates and normalized equation}

Let \(y\to+\infty\), put
\begin{equation}
 Y=\log y,
 \qquad
 w=W_0(4Y\e^{-3}),
 \qquad
 Z=\e^{(w+3)/2}=2\sqrt{\frac{Y}{w}},
 \label{eq:barnes-coordinates}
\end{equation}
and define
\begin{equation}
 D=w+1,
 \qquad s=\log Z=\frac{D+2}{2},
 \qquad q=Z^{-1}.
 \label{eq:barnes-Ds}
\end{equation}
Then
\begin{equation}
 \frac12Z^2\log Z-\frac34Z^2=Y.
 \label{eq:barnes-leading-exact}
\end{equation}
Write
\begin{equation}
 G_+^{\cinv}(y)=1+Z+U_G(q,D),
 \qquad U_G=\sum_{n\ge0}C_n(D)q^n.
 \label{eq:barnes-U-def}
\end{equation}
Dividing the residual by \(\tfrac12Z\) gives
\begin{align}
 \mathscr E_G(q,U;D)={}&
 \frac{(1+qU)^2\left(\dfrac w2+\log(1+qU)\right)-\dfrac w2}{q}
 +\lambda(1+qU)
 \notag\\
 &+q\left[-\frac16\left(s+\log(1+qU)\right)+2\kappa\right]
 \notag\\
 &+\sum_{k\ge1}2b_kq^{2k+1}(1+qU)^{-2k}=0.
 \label{eq:E-barnes}
\end{align}
The coefficient of \(U\) at \(q=0\) is \(D\), and the linear Barnes term
immediately gives \(C_0=-\lambda/D\).

\subsection{Explicit coefficient functions}

\begin{theorem}[Barnes inverse through \(Z^{-3}\)]\label{thm:barnes-main}
Let \(Y,w,Z,D\) be given by
\cref{eq:barnes-coordinates,eq:barnes-Ds}, and let
\(\lambda=\log(2\pi)\), \(\alpha=\log A\).  Then
\begin{equation}
 \BarnesInv(y)=1+Z+C_0+\frac{C_1}{Z}+\frac{C_2}{Z^2}
 +\frac{C_3}{Z^3}+\Ord(Z^{-4}),
 \label{eq:barnes-main-series}
\end{equation}
where
\begin{align}
 C_0={}&-\frac{\lambda}{D},
 \label{eq:barnes-C0}\\
 C_1={}&\frac1{12}+\frac{2\alpha}{D}
 +\frac{\lambda^2}{2D^2}-\frac{\lambda^2}{D^3},
 \label{eq:barnes-C1}\\
 C_2={}&\frac{4\lambda\alpha}{D^3}
 +\frac{4\lambda^3}{3D^4}-\frac{2\lambda^3}{D^5},
 \label{eq:barnes-C2}\\
 C_3={}&-\frac1{288}
 +\frac{11/720-\alpha/6}{D}
 -\frac{2\alpha^2+\lambda^2/24}{D^2}
 \notag\\
 &-\frac{4\alpha^2+\alpha\lambda^2+\lambda^2/12}{D^3}
 -\frac{2\alpha\lambda^2+\lambda^4/8}{D^4}
 \notag\\
 &+\frac{12\alpha\lambda^2-\lambda^4/6}{D^5}
 +\frac{25\lambda^4}{6D^6}-\frac{5\lambda^4}{D^7}.
 \label{eq:barnes-C3}
\end{align}
\end{theorem}

\begin{proof}
Apply the recurrence \cref{eq:triangular-recurrence} to
\cref{eq:E-barnes}.  The first equation is \(DC_0+\lambda=0\).  At each
subsequent order, all nonlinear terms contain only previously determined
coefficients.  Exact simplification with
\(\kappa=\tfrac1{12}-\alpha\) gives
\cref{eq:barnes-C1,eq:barnes-C2,eq:barnes-C3}.  Substitution into
\cref{eq:E-barnes} cancels its coefficients through order \(q^3\), and
residual promotion gives the real asymptotic statement.
\end{proof}

The next coefficient, useful for numerical work, remains reasonably
compact:
\begin{align}
 C_4={}&\frac{11\lambda}{360D^2}+\frac{11\lambda}{360D^3}
 -\frac{\lambda(72\alpha^2+\lambda^2)}{9D^4}
 \notag\\
 &-\frac{\lambda(72\alpha^2+16\alpha\lambda^2+\lambda^2)}{3D^5}
 -\frac{4\lambda^3(50\alpha+3\lambda^2)}{15D^6}
 \notag\\
 &+\frac{4\lambda^3(30\alpha-\lambda^2)}{3D^7}
 +\frac{14\lambda^5}{D^8}-\frac{14\lambda^5}{D^9}.
 \label{eq:barnes-C4}
\end{align}
Thus \cref{eq:barnes-main-series} can be extended by \(C_4/Z^4\) with
remainder \(\Ord(Z^{-5})\), provided the forward expansion is retained to
the corresponding order.

\subsection{A compact hand recurrence}

For low orders it is useful to avoid re-expanding the full normalized
equation.  Let
\begin{equation}
 a=\frac\lambda2,\qquad b=-\frac1{12},\qquad c=\kappa,
 \qquad d_1=-\frac1{240},\qquad d=s-1=\frac D2.
 \label{eq:barnes-short-symbols}
\end{equation}
Writing \(U=\sum c_nZ^{-n}\), one obtains
\begin{align}
 c_0={}&-\frac{a}{d},
 \label{eq:barnes-hand0}\\
 c_1={}&-\frac{ac_0+bs+c+(s/2)c_0^2}{d},
 \label{eq:barnes-hand1}\\
 c_2={}&-\frac{ac_1+bc_0+c_0^3/6+sc_0c_1}{d},
 \label{eq:barnes-hand2}\\
 c_3={}&-\frac{ac_2-(b/2)c_0^2+bc_1+d_1-c_0^4/24
 +(c_0^2c_1)/2+sc_0c_2+(s/2)c_1^2}{d}.
 \label{eq:barnes-hand3}
\end{align}
These are algebraically equivalent to
\cref{eq:barnes-C0,eq:barnes-C1,eq:barnes-C2,eq:barnes-C3}.

\subsection{Derivative check and Newton refinement}

Let
\begin{equation}
 \Psi_G(x)=\frac{\dd}{\dd x}\log G(x).
 \label{eq:barnes-digamma}
\end{equation}
Then
\begin{equation}
 \frac{\dd}{\dd Y}\BarnesInv(\e^Y)
 =\frac1{\Psi_G(\BarnesInv(\e^Y))}.
 \label{eq:barnes-inverse-derivative}
\end{equation}
A numerical implementation can refine any transseries truncation by
\begin{equation}
 x_{\rm new}=x-\frac{\log G(x)-Y}{\Psi_G(x)}.
 \label{eq:barnes-Newton}
\end{equation}
Since \(\Psi_G(1+Z)\sim Z\log Z\), the asymptotic seed is already deep in
the quadratic convergence basin for large \(Y\).

\section{Nested-logarithm re-expansions and sector ordering}

\subsection{Lambert W at infinity}

For \(\xi\to\infty\), put
\begin{equation}
 L=\log\xi,
 \qquad \ell=\log L.
 \label{eq:Lell}
\end{equation}
The standard expansion of \(W_0(\xi)\) is a polynomial-logarithmic series
in \(L^{-1}\) with coefficients in \(\mathbb Q[\ell]\)
\cite{Corless1996,DLMFLambert}.  Re-expanding the exact Lambert normalizers
therefore gives conventional transseries in nested logarithms.

\subsection{Gamma nested-log sector}

For Gamma take \(\xi=T/\e\), so
\begin{equation}
 L=\log(T/\e),\qquad \ell=\log L.
 \label{eq:gamma-Lell}
\end{equation}
Then
\begin{align}
 X=\frac{T}{L}\bigg[{}&1+\frac{\ell}{L}
 +\frac{\ell(\ell-1)}{L^2}
 +\frac{\ell(\ell-2)(2\ell-1)}{2L^3}
 \notag\\
 &+\frac{\ell(6\ell^3-26\ell^2+27\ell-6)}{6L^4}
 +\Ord\!\left(\frac{\ell^5}{L^5}\right)\bigg].
 \label{eq:gamma-X-nested}
\end{align}
Combining this with \cref{eq:gamma-main-expansion} gives a fully elementary
nested-log transseries.

\begin{cautionbox}[title={Sector ordering matters}]
For every fixed \(N\), the omitted term \(X/L^N\) still tends to infinity.
Therefore the additive shift \(\tfrac12\) is smaller than \emph{every fixed}
number of nested-log corrections to \(X\).  Likewise, the first Stirling
inverse sector \(1/(24sX)\) lies beyond the constant sector.  A correctly
ordered answer either retains \(W\) exactly or treats the completed
Lambert subseries as one block before appending \(\tfrac12\) and the
powers of \(X^{-1}\).
\end{cautionbox}

The asymptotic hierarchy is
\begin{equation}
 X\gg \frac{X}{L}\gg\frac{X}{L^2}\gg\cdots
 \gg1\gg\frac1{X\log X}\gg\frac1{X^3\log X}\gg\cdots.
 \label{eq:gamma-hierarchy}
\end{equation}

\subsection{Barnes nested-log sector}

For Barnes take
\begin{equation}
 L=\log(4Y\e^{-3}),\qquad \ell=\log L.
 \label{eq:barnes-Lell}
\end{equation}
Then
\begin{align}
 Z=\frac{2\sqrt Y}{\sqrt L}\bigg[{}&1+\frac{\ell}{2L}
 +\frac{\ell(3\ell-4)}{8L^2}
 +\frac{\ell(5\ell^2-16\ell+8)}{16L^3}
 \notag\\
 &+\frac{\ell(105\ell^3-568\ell^2+720\ell-192)}{384L^4}
 +\Ord\!\left(\frac{\ell^5}{L^5}\right)\bigg].
 \label{eq:barnes-Z-nested}
\end{align}
The first slow denominator has expansion
\begin{align}
 \frac1{w+1}={}&\frac1L+\frac{\ell-1}{L^2}
 +\frac{\ell^2-3\ell+1}{L^3}
 \notag\\
 &+\frac{\ell^3-\tfrac{11}{2}\ell^2+6\ell-1}{L^4}
 +\Ord\!\left(\frac{\ell^4}{L^5}\right).
 \label{eq:inv-D-nested}
\end{align}
Thus, after the completed Lambert block, the Barnes hierarchy begins
\begin{equation}
 Z\quad\hbox{(with all nested logs)},\qquad
 1,\qquad -\frac{\lambda}{w+1},\qquad
 \frac{C_1}{Z},\qquad\frac{C_2}{Z^2},\ldots.
 \label{eq:barnes-hierarchy}
\end{equation}
Again, retaining \(W\) exactly is both cleaner and numerically superior.

\subsection{Complex branches}

Let \(Y_j=\Log y+2\pi\mathrm i j\).  For Gamma, define
\begin{equation}
 T_j=Y_j-c_\Gamma,
 \qquad w_{j,k}=W_k(T_j/\e),
 \qquad X_{j,k}=T_j/w_{j,k}.
 \label{eq:gamma-complex-branches}
\end{equation}
For Barnes, define
\begin{equation}
 w_{j,k}=W_k(4Y_j\e^{-3}),
 \qquad Z_{j,k}=2\sqrt{Y_j/w_{j,k}},
 \label{eq:barnes-complex-branches}
\end{equation}
with a compatible square-root branch.  In sectors where the relevant
Stirling expansions hold and the selected forward branch is univalent,
the same rational coefficient functions give formal inverse sheets.  This
construction does not by itself resolve global branch cuts; it is a local
asymptotic parametrization of the sheets at infinity.  Global inverse
Gamma branch geometry is discussed by Jeffrey and Watt~\cite{JeffreyWatt}.

\section{Numerical verification}

All values in this section were computed with 80 decimal digits.  The
``exact'' inverse is the high-precision root of
\(\log\Gamma(x)=Y\) or \(\log G(x)=Y\).  The tables report absolute error.
They are intended as coefficient checks, not as optimized numerical
algorithms.

\subsection{Gamma}

Let
\begin{equation}
 \mathcal A_{2N-1}(Y)=\frac12+X+
 \sum_{n=1}^{N}\frac{A_{2n-1}(s)}{X^{2n-1}},
 \qquad \mathcal A_{-1}=\frac12+X.
 \label{eq:A-trunc}
\end{equation}

\begin{table}[htbp]
\centering
\caption{Absolute error \(\lvert\mathcal A_m(Y)-\Gamma_+^{\cinv}(\e^Y)\rvert\).}
\label{tab:gamma-errors}
\scriptsize
\setlength{\tabcolsep}{3.2pt}
\begin{tabular}{rcccccc}
\toprule
\(Y\) & \(m=-1\) & \(m=1\) & \(m=3\) & \(m=5\) & \(m=7\) & \(m=9\)\\
\midrule
10
& \(2.4064\!\times10^{-3}\) & \(2.9453\!\times10^{-6}\)
& \(1.2976\!\times10^{-8}\) & \(1.3290\!\times10^{-10}\)
& \(2.5991\!\times10^{-12}\) & \(8.3584\!\times10^{-14}\)\\
50
& \(5.7806\!\times10^{-4}\) & \(8.0628\!\times10^{-8}\)
& \(4.6010\!\times10^{-11}\) & \(6.2767\!\times10^{-14}\)
& \(1.6224\!\times10^{-16}\) & \(6.8367\!\times10^{-19}\)\\
100
& \(3.0464\!\times10^{-4}\) & \(1.5313\!\times10^{-8}\)
& \(3.2682\!\times10^{-12}\) & \(1.6759\!\times10^{-15}\)
& \(1.6235\!\times10^{-18}\) & \(2.5579\!\times10^{-21}\)\\
1000
& \(3.4011\!\times10^{-5}\) & \(4.4431\!\times10^{-11}\)
& \(2.6678\!\times10^{-16}\) & \(3.8769\!\times10^{-21}\)
& \(1.0573\!\times10^{-25}\) & \(4.6656\!\times10^{-30}\)\\
\bottomrule
\end{tabular}
\end{table}

The corresponding exact values begin
\begin{align*}
 \Gamma_+^{\cinv}(\e^{10})&=8.715310138772789660993299\ldots,\\
 \Gamma_+^{\cinv}(\e^{50})&=23.48931723598978590000318\ldots,\\
 \Gamma_+^{\cinv}(\e^{100})&=38.18456559880504784450257\ldots,\\
 \Gamma_+^{\cinv}(\e^{1000})&=226.5074413075941080947948\ldots.
\end{align*}

\subsection{Barnes G}

Define
\begin{equation}
 \mathcal B_N(Y)=1+Z+\sum_{n=0}^{N}\frac{C_n(D)}{Z^n},
 \qquad \mathcal B_{-1}=1+Z.
 \label{eq:B-trunc}
\end{equation}

\begin{table}[htbp]
\centering
\caption{Absolute error \(\lvert\mathcal B_N(Y)-G_+^{\cinv}(\e^Y)\rvert\).}
\label{tab:barnes-errors}
\scriptsize
\setlength{\tabcolsep}{3.0pt}
\begin{tabular}{rcccccc}
\toprule
\(Y\) & \(N=-1\) & \(N=0\) & \(N=1\) & \(N=2\) & \(N=3\) & \(N=4\)\\
\midrule
10
& \(9.3803\!\times10^{-1}\) & \(5.5074\!\times10^{-2}\)
& \(9.5222\!\times10^{-3}\) & \(5.4316\!\times10^{-4}\)
& \(1.3042\!\times10^{-4}\) & \(3.7180\!\times10^{-5}\)\\
50
& \(6.3838\!\times10^{-1}\) & \(3.1685\!\times10^{-2}\)
& \(1.2981\!\times10^{-3}\) & \(5.1633\!\times10^{-5}\)
& \(7.7735\!\times10^{-6}\) & \(1.8663\!\times10^{-8}\)\\
100
& \(5.5115\!\times10^{-1}\) & \(2.2811\!\times10^{-2}\)
& \(5.1226\!\times10^{-4}\) & \(2.4782\!\times10^{-5}\)
& \(1.5720\!\times10^{-6}\) & \(4.4894\!\times10^{-8}\)\\
1000
& \(3.6598\!\times10^{-1}\) & \(7.0937\!\times10^{-3}\)
& \(2.3648\!\times10^{-5}\) & \(9.6114\!\times10^{-7}\)
& \(5.4021\!\times10^{-9}\) & \(2.6721\!\times10^{-10}\)\\
10000
& \(2.6677\!\times10^{-1}\) & \(2.1957\!\times10^{-3}\)
& \(1.2367\!\times10^{-6}\) & \(3.1269\!\times10^{-8}\)
& \(9.7000\!\times10^{-12}\) & \(8.4220\!\times10^{-13}\)\\
\bottomrule
\end{tabular}
\end{table}

The exact roots begin
\begin{align*}
 G_+^{\cinv}(\e^{10})&=6.919323766771216637535535\ldots,\\
 G_+^{\cinv}(\e^{50})&=11.07407747566167840724942\ldots,\\
 G_+^{\cinv}(\e^{100})&=13.92651070973114185567169\ldots,\\
 G_+^{\cinv}(\e^{1000})&=32.55204340388853580099282\ldots,\\
 G_+^{\cinv}(\e^{10000})&=83.54283578073789743023524\ldots.
\end{align*}

At \(Y=10\), the Barnes scale is only moderately large and the asymptotic
character is visible.  By \(Y=50\), the \(C_4\)-truncation already gains
about eight digits, and the improvement accelerates rapidly.

\section{Algorithms and symbolic implementation}

\subsection{Coefficient-generation algorithm}

For a forward expansion of the form \cref{eq:general-forward}, the direct
algorithm is:
\begin{enumerate}
\item choose the leading Lambert coordinate \(w,X,q\);
\item build the normalized residual \(\mathscr E(q,U;w)\);
\item initialize \(U=0\);
\item at order \(n\), append an unknown \(c_nq^n\), extract
      \(\coeff qn{\mathscr E}\), solve the resulting linear equation, and
      simplify;
\item substitute the truncation back into \(\mathscr E\) to verify the
      residual order.
\end{enumerate}
Only forward coefficients through outer order \(n+1\) are needed for
\(C_n\).  The method works over exact rational-function fields, so no
floating-point recognition is necessary.

\subsection{Generic SymPy skeleton}

The following compact code implements \cref{eq:triangular-recurrence} for
\(P_m(s)\) supplied as exact SymPy expressions.  It is deliberately
minimal; production code should cache truncated powers and use Newton
series iteration at high order.

\begin{lstlisting}[style=pythonstyle]
import sympy as sp


def lambert_inverse_coefficients(p, w, s, R, order):
    """Return C_0,...,C_order for the normalized equation (4.5).

    R is a dict {m: R_m(s)}.  Missing m are interpreted as zero.
    The relation s = beta + w/p is imposed by the caller.
    """
    q = sp.Symbol("q")
    coeffs = []
    U = sp.Integer(0)

    for n in range(order + 1):
        cn = sp.Symbol(f"c{n}")
        trial = U + cn*q**n
        h = sp.log(1 + q*trial)

        E = ((1 + q*trial)**p * (w/p + h) - w/p) / q
        for m in range(1, order + 2):
            Rm = R.get(m, sp.Integer(0))
            if Rm != 0:
                shifted = Rm.subs(s, s + h)
                E += q**(m - 1) * (1 + q*trial)**(p - m) * shifted

        eqn = sp.series(E, q, 0, n + 1).removeO().expand().coeff(q, n)
        value = sp.solve(sp.Eq(eqn, 0), cn)[0]
        value = sp.factor(sp.cancel(value))
        coeffs.append(value)
        U += value*q**n

    return coeffs
\end{lstlisting}

For Gamma, use \(p=1\), \(w=s-1\), and
\(R_{2k}=a_k\).  For Barnes, use \(p=2\),
\(w=2s-3\),
\begin{equation}
 R_1=\lambda,
 \qquad R_2=-\frac{s}{6}+2\kappa,
 \qquad R_{2k+2}=2b_k.
 \label{eq:R-concrete}
\end{equation}
The output reproduces every coefficient displayed in this article.

\subsection{Independent verification protocol}

A robust symbolic workflow performs three checks:
\begin{enumerate}
\item \emph{Residual check:} substitute the truncation into the normalized
      equation and verify that all targeted powers of \(q\) vanish.
\item \emph{Derivative check:} differentiate the truncation and compare
      with the exact inverse-derivative identity.
\item \emph{Numerical check:} evaluate at several separated target sizes
      and compare with high-precision root finding.
\end{enumerate}
Agreement of all three checks makes transcription or branch errors much
less likely.

\section{Extensions and structural consequences}

\subsection{More general slow coefficient fields}

Nothing in \cref{thm:formal-inverse} requires \(P_m\) to be a polynomial in
one logarithm.  One may take a slow field generated by
\begin{equation}
 \ell_1=\log X,
 \qquad \ell_{j+1}=\log\ell_j,
 \label{eq:iterated-logs}
\end{equation}
closed under the slow derivation
\begin{equation}
 \delta=
 \frac{\partial}{\partial\ell_1}
 +\frac1{\ell_1}\frac{\partial}{\partial\ell_2}
 +\frac1{\ell_1\ell_2}\frac{\partial}{\partial\ell_3}+\cdots.
 \label{eq:slow-derivation}
\end{equation}
Then \(\dd A/\dd X=q\delta A\), so differentiation raises the primary
valuation exactly as required.  Hahn-field and transseries foundations for
such constructions are developed by van der Hoeven~\cite{vanDerHoeven}.

\subsection{Regular variation viewpoint}

The leading maps \(x^p\log x\) are regularly varying with index \(p\).
General regular-variation theory guarantees asymptotic inverses and
organizes slowly varying factors~\cite{BGT}.  The Lambert normalization is
a stronger, application-specific statement: it inverts the dominant slowly
varying factor exactly and exposes an effective coefficient calculus for
all lower blocks.

\subsection{Multiple logarithmic factors}

For leaders such as
\begin{equation}
 a x^p(\log x)^\rho(\log\log x)^\sigma,
 \label{eq:multiple-log-leader}
\end{equation}
there is generally no one-step elementary Lambert reduction.  A useful
strategy is to invert the first two factors by \cref{eq:rho-normalizer},
treat \((\log\log x)^\sigma\) as a slow multiplier, and solve the remaining
slow equation by fixed-point iteration.  Once a leading inverse \(X\) is
known to relative error \(o(1)\), the additive correction still lives in a
\(q=X^{-1}\) transseries and the triangular residual method applies.

\subsection{Parameter-dependent families}

If \(a,p,\beta\), or the coefficients \(P_m\) depend analytically on a
parameter, all recurrences remain valid over the enlarged coefficient
field.  Uniformity requires keeping \(w+1\) away from zero and controlling
forward remainders uniformly.  This is useful for generalized Gamma and
multiple-Gamma families.

\subsection{Coefficient growth and optimal truncation}

The first contribution of a newly appearing forward coefficient \(P_m\)
to the inverse is its negative divided by the normalized derivative
\(w+1\).  Consequently, factorial growth of forward Stirling coefficients
is inherited by the inverse coefficients.  For Gamma,
\begin{equation}
 a_n\sim
 \frac{2(-1)^n(2n-2)!}{(2\pi)^{2n}},
 \qquad
 A_{2n-1}(s)\sim-\frac{a_n}{s},
 \label{eq:gamma-growth}
\end{equation}
up to the conventional sign interpretation of Bernoulli asymptotics.
This predicts least-term truncation near \(n\asymp\pi X\).  The analogous
Barnes estimate follows from \(b_n\), with least term near
\(n\asymp\pi Z\).  These estimates agree with the exponential scales in
\cref{eq:gamma-exp-scale,eq:barnes-exp-scale}.

\section{Conclusions}

The large-value inverse Gamma and inverse Barnes \(G\) functions share a
single structural mechanism.  Their logarithms have a dominant
power--logarithmic block, and that block is exactly invertible by Lambert
\(W\).  In the adapted coordinate, the remaining inversion problem is a
near-identity formal equation over a slow coefficient field.  Its
Jacobian is \(w+1\), its coefficient equations are triangular, and its
analytic error is controlled directly by the substitution residual.

For Gamma, the half shift converts Stirling's series into an odd inverse
power series and yields an inverse expansion containing only
\(X^{-1},X^{-3},X^{-5},\ldots\) beyond the leading block.  The separate pole
branches are ordinary convergent inversions of the reciprocal Gamma
function.  For Barnes \(G\), the same calculus begins with a nonzero slow
correction \(-\log(2\pi)/(w+1)\) and then produces explicit powers of
\(Z^{-1}\).

The resulting method is not tied to either special function.  It applies
to broad families with leaders \(a x^p(\log x-\beta)^\rho\), admits
iterated-log coefficient fields, supports Lagrange and Newton formulations,
and transports exponentially small sectors in a universal way.  The exact
Lambert representation should normally be retained as the primary answer;
nested-log expansion is a secondary re-expression whose internal sector
must be completed before smaller additive corrections are ordered.

\appendix

\section{Coefficient derivation details}

\subsection{Expansion of the dominant normalized block}

Let
\begin{equation}
 \Phi(t;w)=(1+t)^p\left(\frac wp+\log(1+t)\right).
 \label{eq:Phi-dominant}
\end{equation}
Then
\begin{equation}
 \Phi(t;w)-\Phi(0;w)
 =(w+1)t+\Lambda_2t^2+\Lambda_3t^3+\cdots,
 \label{eq:Phi-expand}
\end{equation}
with \(\Lambda_2,\Lambda_3\) as in
\cref{eq:Lambda2,eq:Lambda3}.  Since \(t=qU\), division by \(q\) produces
\begin{equation}
 (w+1)U+\Lambda_2qU^2+\Lambda_3q^2U^3+\cdots.
 \label{eq:dominant-U-expand}
\end{equation}
Similarly,
\begin{align}
 &(1+t)^{p-1}R_1(s+\log(1+t))
 \notag\\
 &\qquad=r_1+B_1t+B_2t^2+\Ord(t^3),
 \label{eq:R1-expand}\\
 &(1+t)^{p-2}R_2(s+\log(1+t))
 =r_2+D_2t+\Ord(t^2).
 \label{eq:R2-expand}
\end{align}
Collecting powers of \(q\) in
\cref{eq:dominant-U-expand,eq:R1-expand,eq:R2-expand} gives the universal
coefficients \cref{eq:C0-univ,eq:C1-univ,eq:C2-univ}.

\subsection{Why the Barnes constants simplify}

DLMF gives an asymptotic form involving \(z\log\Gamma(z+1)\) and
Glaisher's constant.  Substitution of Stirling's series yields a constant
\(1/12-\log A\), which is \(\zeta'(-1)\).  The two contributions to
\(z^{-2k}\) combine as
\begin{align}
 \frac{B_{2k+2}}{(2k+2)(2k+1)}
 +\frac{B_{2k+2}}{2k(2k+1)(2k+2)}
 =\frac{B_{2k+2}}{2k(2k+2)}.
 \label{eq:barnes-coeff-combine}
\end{align}
This proves \cref{eq:barnes-forward,eq:bk-def} from the DLMF form.

\section{Further pole-branch coefficients}

Let \(t=(-1)^n/(n!y)\), \(\psi_n=H_n-\gamma\), and
\(H_n^{(r)}=\sum_{k=1}^nk^{-r}\).  Continuing
\cref{eq:pole-Lagrange} gives
\begin{align}
 \Gamma_{[-n]}^{\cinv}(y)={}&-n+t+\psi_nt^2
 +\frac{3\psi_n^2+H_n^{(2)}+\zeta(2)}2t^3
 \notag\\
 &+\frac{8\psi_n^3+6\psi_n(H_n^{(2)}+\zeta(2))
 +H_n^{(3)}-\zeta(3)}3t^4
 +\Ord(t^5).
 \label{eq:pole-four-terms}
\end{align}
For \(n=0\), this reduces to the first four terms of
\cref{eq:gamma-left-expansion}.

\section{Notation and reproducibility checklist}

\begin{longtable}{@{}p{0.22\textwidth}p{0.71\textwidth}@{}}
\toprule
Symbol & Meaning\\
\midrule
\endfirsthead
\toprule
Symbol & Meaning\\
\midrule
\endhead
\(f^{\cinv}\) & Compositional inverse, not reciprocal.\\
\(Y\) & Logarithm of the large function value: \(Y=\log y\).\\
\(w\) & Lambert coordinate that exactly inverts the dominant block.\\
\(X,Z\) & Exact leading inverses for Gamma and Barnes \(G\), respectively.\\
\(q\) & Primary small parameter, \(q=X^{-1}\) or \(Z^{-1}\).\\
\(s,D\) & Slow variables: \(s=w+1\) for Gamma and \(D=w+1\) for Barnes.\\
\(A\) & Glaisher--Kinkelin constant; \(\alpha=\log A\).\\
\(\mathscr E\) & Normalized substitution residual.\\
\bottomrule
\end{longtable}

To reproduce the results:
\begin{enumerate}
\item generate Bernoulli or Bernoulli-polynomial coefficients exactly;
\item use \cref{eq:E-gamma} or \cref{eq:E-barnes};
\item apply \cref{eq:triangular-recurrence};
\item check cancellation in the normalized residual;
\item evaluate against high-precision roots of the logarithmic equations.
\end{enumerate}

\begin{thebibliography}{99}

\bibitem{BGT}
N.~H. Bingham, C.~M. Goldie, and J.~L. Teugels,
\newblock \emph{Regular Variation},
\newblock Encyclopedia of Mathematics and its Applications 27,
Cambridge University Press, Cambridge, 1987.

\bibitem{BorweinCorless}
J.~M. Borwein and R.~M. Corless,
\newblock Gamma and factorial in the Monthly,
\newblock \emph{American Mathematical Monthly} \textbf{125} (2018),
400--424,
\newblock DOI: \href{https://doi.org/10.1080/00029890.2018.1420983}
{10.1080/00029890.2018.1420983}.

\bibitem{Comtet}
L.~Comtet,
\newblock \emph{Advanced Combinatorics: The Art of Finite and Infinite
Expansions},
\newblock D.~Reidel, Dordrecht, 1974.

\bibitem{Corless1996}
R.~M. Corless, G.~H. Gonnet, D.~E.~G. Hare, D.~J. Jeffrey, and D.~E. Knuth,
\newblock On the Lambert \(W\) function,
\newblock \emph{Advances in Computational Mathematics} \textbf{5} (1996),
329--359.

\bibitem{DLMFLambert}
NIST Digital Library of Mathematical Functions,
\newblock \S4.13, Lambert \(W\)-function,
\newblock F.~W.~J. Olver et al., editors,
\newblock \url{https://dlmf.nist.gov/4.13}.

\bibitem{DLMFGamma}
NIST Digital Library of Mathematical Functions,
\newblock \S5.11, asymptotic expansions of the Gamma function,
\newblock F.~W.~J. Olver et al., editors,
\newblock \url{https://dlmf.nist.gov/5.11}.

\bibitem{DLMFBarnes}
NIST Digital Library of Mathematical Functions,
\newblock \S5.17, Barnes' \(G\)-function,
\newblock F.~W.~J. Olver et al., editors,
\newblock \url{https://dlmf.nist.gov/5.17}.

\bibitem{FlajoletSedgewick}
P.~Flajolet and R.~Sedgewick,
\newblock \emph{Analytic Combinatorics},
\newblock Cambridge University Press, Cambridge, 2009.

\bibitem{Henrici}
P.~Henrici,
\newblock \emph{Applied and Computational Complex Analysis}, Vol.~2,
\newblock Wiley, New York, 1977.

\bibitem{JeffreyWatt}
D.~J. Jeffrey and S.~M. Watt,
\newblock The inverse of the complex Gamma function,
\newblock arXiv:2311.16583, 2023,
\newblock \url{https://arxiv.org/abs/2311.16583}.

\bibitem{Nemes}
G.~Nemes,
\newblock Error bounds and exponential improvement for the asymptotic
expansion of the Barnes \(G\)-function,
\newblock \emph{Proceedings of the Royal Society A} \textbf{470} (2014),
20140534,
\newblock DOI: \href{https://doi.org/10.1098/rspa.2014.0534}
{10.1098/rspa.2014.0534}.

\bibitem{Olver}
F.~W.~J. Olver,
\newblock \emph{Asymptotics and Special Functions},
\newblock A K Peters, Wellesley, MA, 1997.

\bibitem{ParisKaminski}
R.~B. Paris and D.~Kaminski,
\newblock \emph{Asymptotics and Mellin--Barnes Integrals},
\newblock Cambridge University Press, Cambridge, 2001.

\bibitem{Pedersen}
H.~L. Pedersen,
\newblock Inverses of Gamma functions,
\newblock \emph{Constructive Approximation} \textbf{41} (2015), 251--267,
\newblock DOI: \href{https://doi.org/10.1007/s00365-014-9239-1}
{10.1007/s00365-014-9239-1}; arXiv:1309.2167.

\bibitem{ProveItTransseries}
V.~Reshetnikov,
\newblock Series and transseries draft packages,
\newblock ProveIt repository,
\newblock \href{https://github.com/VladimirReshetnikov/ProveIt/tree/main/Analysis/FabiusFunction/docs/semi-formalized-research-frontiers/drafts/series-and-transseries}
{repository directory \texttt{series-and-transseries}}.

\bibitem{Uchiyama}
M.~Uchiyama,
\newblock The principal inverse of the Gamma function,
\newblock \emph{Proceedings of the American Mathematical Society}
\textbf{140} (2012), 1343--1348,
\newblock DOI: \href{https://doi.org/10.1090/S0002-9939-2011-11023-2}
{10.1090/S0002-9939-2011-11023-2}.

\bibitem{vanDerHoeven}
J.~van der Hoeven,
\newblock \emph{Transseries and Real Differential Algebra},
\newblock Lecture Notes in Mathematics 1888, Springer, Berlin, 2006.

\end{thebibliography}

\end{document}
